\section{Support, Locality and Transitive Exposure}
\label{sec:theory}

\subsection{Support Through Admissible Witnesses}

For a proposition $c$, let $\mathcal W(c)$ be a finite family of admissible
witnesses. Each $W$ is a set of jointly required evidence or prerequisite
propositions; different $W$ are alternatives. We use
\begin{equation}
 \operatorname{Supported}(c,t;q)=
 \bigvee_{W\in\mathcal W(c)}
 \left[\psi_c(W,t;q)\land
       \bigwedge_{u\in W}\operatorname{Valid}(u,t;q)\right].
\label{eq:support}
\end{equation}
$\operatorname{Valid}$ requires observable availability and admissible temporal
and subject scope; for a prerequisite claim it recursively requires support.
$\psi_c$ checks the asserted interval, selected logical versions, quantity,
normalized unit and numerical operation. Numerical agreement uses the frozen
tolerance $|v-\widehat v|\leq\max(10^{-6},0.01|\widehat v|)$.
A correction can therefore invalidate $\psi_c$ while leaving the source
available. A supported historical assertion at its original knowledge time and
a supported current assertion are different propositions.
The extension's specified inequalities use exact floating-point comparisons.

The original model supports local AND/OR composition and a two-operand
difference. Generated numerical claims require all declared parents plus an
admissible direct numerical witness; comparisons cite both operands. The
extension uses explicit OR lists of AND inputs composed through an acyclic
positive claim graph, not unrestricted Boolean formulas. Missing-evidence
statements require a bounded producer universe and explicit availability watches;
absence in an open graph is insufficient.
The extension stores witness grouping in each record's semantic-program metadata;
source IDs contain their union for traversal. Edge presence alone does not
recover the AND/OR grouping.

As in provenance and truth maintenance~\cite{green2007,doyle1979}, losing one
route preserves support when another admissible witness survives. Our controlled
alternatives use disjoint recording windows, not renamed copies of one
measurement. Source independence does not imply statistical independence.

\subsection{Locality of a Revision}

Let $U_\delta$ contain revised input identities, all earlier versions that can
still be cited, and any changed availability subscriptions. In the dependency
projection, edges point from dependent to input. Define
\begin{equation}
 R_\delta=U_\delta\cup\operatorname{RevReach}(U_\delta),\qquad
 B_\delta=\{c:S_{t^-}(c)\ne S_{t^+}(c)\},
 \label{eq:locality}
\end{equation}
where $S_t(c)$ abbreviates Equation~\ref{eq:support}. $R_\delta$ is a candidate
set; $B_\delta$ is the actual support-change set.

\begin{proposition}[Conditional locality]
For a finite acyclic dependency model, $B_\delta\subseteq R_\delta$ when every
mutable input to each support predicate is explicitly represented and only
$U_\delta$ changes between the two snapshots.
\end{proposition}
\noindent\emph{Proof sketch.}
Order nodes from inputs to dependents. A node outside $R_\delta$ has no changed
input and no prerequisite in $R_\delta$; otherwise it would be reverse reachable.
Its scope, numerical operands and availability state are therefore unchanged.
Induction in the acyclic order gives equal prerequisite support and equal
predicate values, hence equal support. Considering all earlier versions covers
immutable citations to $v_1$ after $v_2$ or $v_3$ arrives. Explicit availability
records cover an arrival that falsifies an absence statement.\hfill$\square$

Locality fails for unrecorded clock expiry, newly discovered unsubscribed
witnesses or external predicate changes. Our fixed-target queries and controlled
witness universes represent changes through version/availability events. Benign
updates and surviving alternatives demonstrate why reachability is only a
candidate set, with potentially strict inclusion $B_\delta\subset R_\delta$.

Indexed adjacency traversal costs $O(|V_{R_\delta}|+|E_{R_\delta}|)$, plus current
version lookups and witness evaluation. Actual costs also include snapshots,
serialization, database round trips and writes. The original runtime repeatedly
constructs version maps; structural replay evaluates the whole small program
before scheduled writes. Neo4j variable-length queries may enumerate paths;
SQLite uses a recursive CTE with duplicate elimination. Neither implementation
inherits a linear-time guarantee from the abstract traversal bound.

\subsection{When Direct Maintenance Suffices}

Let $A_\delta$ be the independently determined set of prior admitted claims requiring
a state change, and $D_\delta$ the claims actually scheduled by direct-only
maintenance. For $|A_\delta|>0$, define
\begin{equation}
 \xi_\delta=\frac{|A_\delta\setminus D_\delta|}{|A_\delta|}.
 \label{eq:exposure}
\end{equation}
We report NA when $A_\delta$ is empty. The primary measured exposure concerns
withdrawal of initially grounded claims; restoration is recorded separately as
a support transition. $A_\delta$ comes from the event-log/task oracle, not from
an accepted flag or the existence of a graph path.

\begin{proposition}[Direct/full equivalence and misses]
Assume initially correct states, complete re-evaluation of $D_\delta$, dependency
semantics consistent with the oracle, and no unrelated refresh or other mechanism
that revisits excluded claims. Before fresh generation, direct maintenance misses
the fraction $\xi_\delta$ of required changes. A sufficient condition for direct
and full maintenance to agree is that every changed support predicate is directly
scheduled. In particular, this holds when each changed input affecting an answer
has a direct dependency edge to every affected claim and all other predicates
remain unchanged.
\end{proposition}
\noindent\emph{Proof sketch.}
Complete direct re-evaluation changes every required member of $D_\delta$.
Excluded members retain their initial state by the no-refresh assumption, so the
missed set is exactly $A_\delta\setminus D_\delta$. If all changed predicates
are scheduled, both algorithms compute the same predicates from the same snapshot;
the remaining predicates retain equal states.\hfill$\square$

A nontrivial counterexample uses $c_1=$ ``test 1 holds'' and $c_2=$ ``tests 1 and
2 hold'', with $c_2$ supported by $c_1$ and test 2. Correcting test 1 to false
requires changing both claims. A direct scheduler reaches only $c_1$, giving
$\xi_\delta=1/2$. If $c_2$ also directly cites test 1, it is directly scheduled
and exposure becomes zero. If the answer has another independently sufficient
route, losing the first route need not change the answer at all. Thus path length
alone predicts neither a miss nor a withdrawal. The waveform exports and the
controlled structural experiment test these distinct regimes.
